% \textsurd (√) in Latin Modern, where tex4ht's lm-ts1.htf gives a middle
% dot: in each font of the text, in a footnote, a theorem, a heading and its
% contents entry with hyperref, a list's label, a table and its caption,
% beside \textbullet, in \text in math, and in a TikZ picture, which keeps
% the font's glyph.
\documentclass{article}
\usepackage[T1]{fontenc}
\usepackage{lmodern}
\usepackage{textcomp}
\usepackage{amsmath}
\usepackage{amsthm}
\usepackage{tikz}
\usepackage{hyperref}
\newtheorem{theorem}{Theorem}
\begin{document}
\tableofcontents
\section{Roots \textsurd{} here}
Plain \textsurd{} sign.
\textit{Italic \textsurd{} sign.}
\textbf{Bold \textsurd{} sign.}
\textbf{\textit{Bolditalic \textsurd{} sign.}}
\textsl{Slanted \textsurd{} sign.}
\textsc{Smallcaps \textsurd{} sign.}
\texttt{Typewriter \textsurd{} sign.}
\textsf{Sans \textsurd{} sign.}
\textsf{\textbf{Sansbold \textsurd{} sign.}}
{\large Large \textsurd{} sign.}
Bullet \textbullet{} and surd \textsurd{} together.
Note\footnote{Footnote \textsurd{} sign.} here.
\begin{theorem}
Theorem \textsurd{} sign.
\end{theorem}
\begin{itemize}
\item[\textsurd] Checked item.
\end{itemize}
\begin{table}[h]
\centering
\begin{tabular}{ll}
Cell & \textsurd \\
\end{tabular}
\caption{Caption \textsurd{} sign.}
\end{table}
Math $x = \text{Mathtext \textsurd{} sign}$ here.
Picture \tikz \node {\textsurd}; here.
\end{document}
